% Replacement for Sections 3.2 and 3.3 of the supplied ARR draft.
% Uses only macros/packages already present in that draft.
% Academic revision: main text focuses on the method and its rationale.

\subsection{Adaptive non-target support}
\label{sec:council}

The council should capture complementary predictions among alternatives that
its members consider likely. We therefore define a shared non-target support
at each reasoning position, adapting its size to the predictions of individual
experts. Taking the union of their candidate sets preserves alternatives
favored by any member and provides a common space in which to compare their
predictions.

At position $t$, let $z^{\downarrow}_{m,t,j}$ and $v_{m,t,j}$ denote the
$j$th-largest logit and its token identity among
$\Vocab\setminus\{y_t\}$ for expert $m$. We search only the top
$K=\min(512,|\Vocab|-1)$ non-target candidates. Both rank and logit axes
are normalized within this window, so the criterion identifies a local knee
in the plausible head region rather than a global vocabulary knee. We use
a locally normalized bounded Kneedle-style criterion
\citep{satopaa2011kneedle}:
\begin{equation}
\begin{gathered}
 x_j=\frac{j-1}{K-1},\quad
 u_{m,t,j}=\frac{z^{\downarrow}_{m,t,j}-z^{\downarrow}_{m,t,K}}
 {z^{\downarrow}_{m,t,1}-z^{\downarrow}_{m,t,K}+\epsilon_{\mathrm{sel}}},\\
 \widehat k_{m,t}=\argmax_{1\le j\le K}
       \left[(1-x_j)-u_{m,t,j}\right],\\
 k_{m,t}=\min\!\left(K,\max(\widehat k_{m,t},8)\right).
\end{gathered}
\label{eq:kneedle}
\end{equation}
Here $\epsilon_{\mathrm{sel}}>0$ ensures a well-defined logit normalization.
The first maximizing rank is chosen when there is a tie. For $K=1$, we set
$\widehat k_{m,t}=k_{m,t}=1$ directly. The search window is the only upper
bound; the lower bound of eight avoids impoverished supports whenever enough
candidates are available. The criterion identifies the largest deviation of
the locally normalized logit curve from the descending diagonal.

The council support is
\begin{equation}
 \Ccal_t=\bigcup_{m=1}^{M}
              \{v_{m,t,j}:1\le j\le k_{m,t}\}.
 \label{eq:stage1-support}
\end{equation}
All experts are evaluated on these common token identities, and the support
is held fixed when differentiating the diversity objective. Excluding the
observed token $y_t$ directs this objective toward alternative continuations
of the reasoning trace.

\subsection{Stage 1: Learning a diverse council}
\label{sec:stage1}

Stage~1 learns a council that combines fidelity to the observed traces with
complementary predictions. We train $M$ independently initialized LoRA
adapters on the same fixed trajectories. Let
$\theta_m=\{A_m^l,B_m^l\}_{l\in\Lcal}$ denote the parameters of expert $m$.
The training rule combines supervised learning, a determinantal objective in
prediction space, and a kernel interaction in effective-weight space.

\paragraph{Trace supervision.}
Each expert learns from the full observed response through supervised
fine-tuning. Let $\Tset_n^{\mathrm{resp}}$ denote its response positions and
$p_{m,n,t}$ the predictive distribution conditioned on the problem and
observed prefix. The likelihood objective is
\begin{equation}
 \Lcal_{\mathrm{SFT}}^{(m)}=
 -\frac{1}{N_{\mathrm{resp}}}
 \sum_{n\in\Bcal}\sum_{t\in\Tset_n^{\mathrm{resp}}}
 \log p_{m,n,t}(y_{n,t}),
 \label{eq:sft1}
\end{equation}
where $N_{\mathrm{resp}}=\sum_{n\in\Bcal}|\Tset_n^{\mathrm{resp}}|$.
This term anchors each expert to the supplied reasoning and answer, while the
council interactions encourage complementary behavior among alternatives.

\paragraph{Diversity in prediction space.}
We measure predictive diversity on the non-target support from
\S\ref{sec:council}. For each reasoning position, define the normalized
feature vector
\begin{equation}
 v_{m,t}(v)=\frac{p_{m,t}(v)}
                  {\sqrt{\sum_{u\in\Ccal_t}p_{m,t}(u)^2}},
 \qquad v\in\Ccal_t.
 \label{eq:stage1-features}
\end{equation}
This representation depends on the relative probabilities of the alternatives,
independently of their total mass. Let $\mathbf V_t$ stack these vectors as
rows. For reasoning step $s$ of example $n$, with content positions
$\Tset_{n,s}^{\mathrm r}$, define
\begin{equation}
 G_{n,s}=\frac{1}{|\Tset_{n,s}^{\mathrm r}|}
           \sum_{t\in\Tset_{n,s}^{\mathrm r}}
                   \mathbf V_t\mathbf V_t^{\top}.
 \label{eq:stage1-gram}
\end{equation}
This is the Gram matrix of the experts' concatenated, normalized step
representations. We encourage these representations to span distinct
directions using a determinantal penalty \citep{kulesza2012dpp}:
\begin{equation}
 \begin{gathered}
 \ell_{n,s}=-\frac{1}{M}\log\det(G_{n,s}+\epsilon_{n,s}\mathbf I),\\
 \Lcal_{\mathrm{DPP}}=
 \frac{1}{|\Bcal_{\mathrm r}|}
 \sum_{n\in\Bcal_{\mathrm r}}\frac{1}{|S_n^{\mathrm r}|}
 \sum_{s\in S_n^{\mathrm r}}\ell_{n,s}.
 \end{gathered}
 \label{eq:dpp}
\end{equation}
Here $\epsilon_{n,s}>0$ is a small diagonal regularizer,
$S_n^{\mathrm r}$ contains the nonempty reasoning steps, and
$\Bcal_{\mathrm r}$ contains the examples with such steps. The determinant
acts on the step-averaged Gram matrix, encouraging complementary predictions
across the reasoning step. Equal weighting across steps and examples limits
the influence of differences in trace length. This objective applies only to
reasoning content; SFT supervises the full response.

\paragraph{Diversity in effective-weight space.}
The LoRA factorization is nonunique: for any invertible $R$,
$(B_m^lR)(R^{-1}A_m^l)=B_m^lA_m^l$. We therefore compare experts through
their effective updates $\Delta W_m^l=(\alpha/r)B_m^lA_m^l$, using
\begin{equation}
 d_{mq}^{2}=\frac{1}{|\Lcal|}\sum_{l\in\Lcal}
 \frac{\|\Delta W_m^l-\Delta W_q^l\|_F^2}
      {d_{\mathrm{out}}^l d_{\mathrm{in}}^l}.
 \label{eq:stage1-distance}
\end{equation}
Normalizing by the matrix size gives each adapted projection equal weight.
The resulting distance is invariant to refactorizations that preserve the
effective update.

An RBF kernel $K_{mq}=\exp(-d_{mq}^{2}/h)$ expresses proximity between
experts, where the positive bandwidth $h$ follows a smoothed median of their
pairwise distances. Inspired by kernel-based particle methods
\citep{liu2016svgd}, we define a similarity potential and a repulsive direction:
\begin{equation}
 \Phi=\frac{1}{\binom{M}{2}}\sum_{m<q}K_{mq},
 \qquad r_m=-\nabla_{\theta_m}\Phi.
 \label{eq:stage1-repulsion}
\end{equation}
With $h$ held fixed, $r_m$ is a local descent direction for pairwise similarity,
providing a second diversity signal alongside the prediction-space objective.

\paragraph{Coupled council update.}
Let $g_m^{\mathrm{SFT}}=\nabla_{\theta_m}\Lcal_{\mathrm{SFT}}^{(m)}$ and
$g_m^{\mathrm{DPP}}=\nabla_{\theta_m}\Lcal_{\mathrm{DPP}}$. The kernel
coordinates gradient exchange across experts:
\begin{equation}
 \begin{gathered}
 \omega_{qm}=\frac{K_{qm}}{\sum_{j=1}^{M}K_{jm}},\\
 a_m=\sum_{q=1}^{M}\omega_{qm}
       \bigl(g_q^{\mathrm{SFT}}+\lambda_{\mathrm{div}}g_q^{\mathrm{DPP}}\bigr).
 \end{gathered}
 \label{eq:stage1-mixing}
\end{equation}
The weights form a convex combination for each receiving expert, favoring
updates from members with similar effective weights. Gradients are combined
in corresponding LoRA factor coordinates. We limit the repulsive direction
relative to this aggregate by setting
\begin{equation}
 \widetilde r_m=
 \min\left\{1,\frac{\|a_m\|_2}{\|r_m\|_2+\delta}\right\}r_m,
 \label{eq:stage1-repulsion-cap}
\end{equation}
where $\delta>0$ is a small stabilizer. Thus
$\|\widetilde r_m\|_2\le\|a_m\|_2$.

A scheduling coefficient $\gamma\in[0,1]$ introduces the council interactions
gradually after an initial period of independent SFT. The resulting
pseudo-gradient is
\begin{equation}
 \begin{aligned}
 g_m={}&(1-\gamma)g_m^{\mathrm{SFT}}\\
       &+\gamma\bigl(a_m-\lambda_{\mathrm{rep}}\widetilde r_m\bigr).
 \end{aligned}
 \label{eq:stage1-update}
\end{equation}
The coefficients $\lambda_{\mathrm{div}}$ and $\lambda_{\mathrm{rep}}$ control
the strengths of predictive diversity and weight-space repulsion. Kernel
weights and the repulsion rescaling are fixed when assembling this direction.
Together, gradient exchange and repulsion balance shared trace supervision
with complementary expert behavior.
